% !TEX root = 00_instructivo.tex
\chapter{Transductores eléctricos y magnéticos: Hall, magnetómetro y galgas}

\section*{Objetivos}
\begin{itemize}
    \item Medir campos magnéticos con sensor Hall (ACS712) y magnetómetro QMC5883L.
    \item Detectar deformación mecánica mediante galgas piezorresistivas.
    \item Integrar los sensores en un circuito sobre PCB de una cara.
\end{itemize}

\section*{Materiales y equipo}
\begin{itemize}
    \item Sensor Hall ACS712 (para corriente/campo).
    \item Magnetómetro triple eje QMC5883L (GY-271).
    \item Galgas extensiométricas (resistencia nominal 120 $\Omega$).
    \item Placa de cobre lisa (para PCB).
    \item Láminas de acrílico y polvo ferroso.
    \item Arduino Uno.
    \item Mini protoboard (solo para pruebas).
    \item Amplificador operacional LM358 (para puente de Wheatstone).
\end{itemize}

\section*{Instrucciones}
\begin{enumerate}
    \item Conecte el ACS712: VCC a 5V, GND a GND, OUT a A0.
    \item Conecte el QMC5883L mediante I2C (A4=SDA, A5=SCL).
    \item Para las galgas: monte un puente de Wheatstone con una galga activa y tres resistencias fijas. Amplifique la salida con LM358 y conéctela a A1.
    \item Coloque polvo ferroso entre acrílicos y acerque un imán. Mida el campo con el magnetómetro (componentes X,Y,Z) y con el Hall.
    \item Aplique peso sobre la galga (placa de cobre) y registre la deformación.
    \item Diseñe y fabrique una PCB (transferencia térmica) para interconectar los sensores sin protoboard.
\end{enumerate}

\section*{Esquemático}
\begin{figure}[htbp]
\centering
\begin{circuitikz}[scale=0.7]
\draw (0,0) node[above]{Arduino};
% ACS712
\draw (1,0.5) node[anchor=east]{5V} -- (2,0.5) node[anchor=west]{VCC};
\draw (1,1) node[anchor=east]{GND} -- (2,1) node[anchor=west]{GND};
\draw (1,1.5) node[anchor=east]{A0} -- (2,1.5) node[anchor=west]{OUT};
% QMC5883L I2C
\draw (1,2.5) node[anchor=east]{A4} -- (3,2.5);
\draw (1,3) node[anchor=east]{A5} -- (3,3);
% Puente de galga
\draw (4,1) node[op amp] (opamp) {};
\draw (opamp.out) -- (5,1) node[anchor=west]{A1};
\end{circuitikz}
\caption{Diagrama simplificado. El puente de galgas requiere alimentación 5V y su salida diferencial va al amplificador.}
\end{figure}

\section*{Preguntas de análisis}
\begin{enumerate}
    \item Compare la sensibilidad del ACS712 (efecto Hall lineal) vs el QMC5883L (magnetómetro triaxial). ¿Cuándo usar cada uno?
    \item ¿Cómo afecta la temperatura a la resistencia de la galga? ¿Cómo se compensa?
    \item Explique la ventaja de usar una PCB sobre protoboard en este tipo de sensores (ruido, impedancia).
    \item ¿Qué fenómeno magnético permite medir corriente sin contacto con el ACS712?
\end{enumerate}

\section*{Bibliografía}
\begin{itemize}
    \item ACS712 Datasheet. Allegro MicroSystems.
    \item QMC5883L Datasheet. QST Corporation.
    \item Galgas extensiométricas – Notas de aplicación Vishay.
\end{itemize}
